% =========================================================
% Proof: Cauchy Sequences Have at Most One Limit
% Source: notes/cauchy/notes-cauchy.tex
% Difficulty: Easy
% =========================================================

\subsection*{Cauchy Sequences Have at Most One Limit}
\label{prf:cauchy-unique-limit}

\begin{remark*}[Return]
$\leftarrow$ \hyperref[prop:cauchy-unique-limit]{Back to Proposition in Notes}
\end{remark*}

\begin{proof}
\Claim Let $(X,d)$ be a metric space.
If $(x_n)$ converges, its limit is unique.

\Given A sequence $(x_n)$ in $(X,d)$ with $x_n \to x$ and $x_n \to y$.

\Goal To show $x = y$.

\Strategy Use the triangle inequality to bound $d(x,y)$ by $d(x, x_n) + d(x_n, y)$.
Both terms go to $0$. Since $d(x,y) \ge 0$ and is constant, conclude $d(x,y) = 0$,
hence $x = y$ by M2.

% -------------------------------------------------------
% YOUR PROOF HERE
% -------------------------------------------------------

\end{proof}

\begin{remark*}[Proof shape]
Squeeze: $0 \le d(x,y) \le d(x,x_n) + d(x_n,y) \to 0$.
\end{remark*}

\begin{remark*}[Dependencies]
\begin{itemize}
  \item Definition~\ref{def:metric}: M2 (identity), M4 (triangle inequality).
  \item Definition~\ref{def:convergence-metric}: $x_n \to x$ and $x_n \to y$.
  \item Squeeze theorem for real sequences (both bounds $\to 0$).
\end{itemize}
\end{remark*}
